\section{Evaluation}
\label{sec:evaluation}

% Target: ~2 two-column pages (~4 single-column equivalent)
% Condense report sections/evaluation.tex. Present QEMU limitations upfront,
% then two demo workloads with results, and architectural cost analysis.

\subsection{Evaluation Platform and Limitations}
\label{sub:platform}

% [~0.5 column]
%   - Platform: QEMU system emulation (v7.2.0), RISC-V ``virt'' machine.
%     RISC-V RV32IMAC with Machine and User modes. FreeRTOS-MPU RISC-V port
%     compiled with configENABLE\_DOMAINS = 1.
%   - Cycle counting: RISC-V `rdcycle`/`rdcycleh` CSR instructions read the
%     `mcycle` counter. QEMU increments this per \emph{instruction} rather
%     than per clock cycle; it does not model pipeline stalls, cache misses,
%     or branch mispredictions.
%   - Implication: the evaluation verifies \emph{constant-time execution
%     paths} (instruction-level determinism) but cannot validate
%     microarchitectural temporal isolation. True hardware evaluation
%     requires silicon with fence.t support or a cycle-accurate simulator
%     (gem5). This is discussed in Section VII.
%   - Metrics: instruction count variance across repeated domain-switch
%     cycles; round-trip latency for cross-domain communication.

\subsection{Blinky Demo: Two Domains}
\label{sub:blinky}

% [~0.75 column]
%   - Setup: Domain 0 contains a sender task at slice 4; Domain 1 contains
%     a receiver task at slice 2. Tasks communicate via a shared queue.
%     Sender sends a counter value; receiver echoes it back. Round trip =
%     sender $\to$ receiver $\to$ sender.
%   - Configuration: 8 total time slices, 1 tick per slice.
%   - Measurement: `rdcycle` captured at sender task entry and after queue
%     receive. Repeated over 1000 iterations.
%   - Results:
%     \begin{itemize}
%       \item Mean round-trip instruction count: 50,000 cycles.
%       \item Variance: $\pm 1$ cycle across all 1000 iterations.
%       \item The $\pm 1$ cycle is attributable to a single conditional
%         branch in the queue-receive path (empty vs.\ non-empty queue).
%     \end{itemize}
%   - Interpretation: the domain scheduler produces deterministic,
%     instruction-count-constant execution across domains. No information
%     about other domains is leaked through instruction count.
%   - [Optional figure: bar chart of round-trip cycle counts with error bars
%     showing $\pm 1$.]

\subsection{Multi-Domain Demo: Four Domains}
\label{sub:multi_domain}

% [~0.75 column]
%   - Setup: 4 domains (Domain 0--3), each allocated 2 time slices. Within
%     each domain, a privileged task creates a higher-priority unprivileged
%     task. The unprivileged task preempts within its domain's slice.
%   - Expected behaviour:
%     \begin{itemize}
%       \item Domain ticks advance 0, 1, 2, 3, 4, 5, 6, 7 cyclically.
%       \item Each domain gets exactly its allocated slices per round.
%       \item Intra-domain: the higher-priority unprivileged task correctly
%         preempts the privileged task within the domain's slice.
%       \item Domain boundaries are respected: no domain's execution bleeds
%         into another's slices.
%     \end{itemize}
%   - Results: the observed domain-tick sequence matches the expected
%     deterministic schedule. Cycle counts per domain per round are
%     constant.
%   - Trade-off: tasks that previously ran with low scheduling latency now
%     experience domain-relative delays. A task delayed at slice $N$ must
%     wait until its domain's next slice before resuming. This is the cost
%     of temporal isolation.

\subsection{Architectural Complexity and Throughput Cost}
\label{sub:cost}

% [~0.5 column]
%   - Code complexity: the domain scheduler adds approximately 500 lines of
%     C across the FreeRTOS kernel. The main changes are concentrated in
%     tasks.c (scheduler) and port.c (tick handling).
%   - Throughput impact: domain scheduling introduces mandatory idle time.
%     When a domain has no ready tasks, its time slices are wasted. In the
%     worst case (one active domain, $N$ domains total), throughput is
%     reduced to $1/N$ of unmodified FreeRTOS.
%   - `vTaskDelayUntil()` semantics change: a task that calls
%     `vTaskDelayUntil()` may oversleep if the target wake time falls in
%     another domain's slice. This breaks the real-time guarantee of the
%     existing API.
%   - Memory overhead: state arrays scale linearly with
%     `configNUM\_TIME\_SLICES` $\times$ `configMAX\_PRIORITIES`. Configurable
%     at compile time.
